\section{The symplectic form and mean evolution on time blocks}
\label{S:actual-symplectic-phase}

Retain the original model, covariance amplitude, physical mask $m_R$ and real $E
= H^{1/2}\times H^{-1/2}$. Fix finite $R,T > 0$ and admissible exponents as
in~\eqref{E:driver-parameters}. Let $M = M_R(W)$ be the from the given noise
first interaction path on $E$, and $M^\rho$ its regularized paths from
Proposition~\ref{P:source-gaussian-transfer} and the source construction. Here
$\rho$ is the positive heat/Laplace smoothing pair. The construction preserving
the given noise and ordinary rows are those of
Lemma~\ref{L:source-retaining-isometry},
Lemma~\ref{L:source-pointfield-regularity} and
Proposition~\ref{P:source-ordinary-response}. The full $F$ operator path and
Young evolution are fixed by Proposition~\ref{P:constant-phase-driver} and
Proposition~\ref{P:constant-young-linear}; their $E$ restrictions are used here.
Complexifications commute with conjugation, so these real restrictions are the
physical phases. For fixed $\Phi$ in the full positive-time original $H_W$ and
real $c$ set
\begin{equation}\label{E:symplectic-combined-driver}
  \begin{aligned}
    Z&=W(\Phi),\qquad D_\Phi=\E[ZM],\qquad
    D_\Phi^\rho=\E[ZM^\rho],\\
    D&=cM+D_\Phi,\qquad D^\rho=cM^\rho+D_\Phi^\rho .
  \end{aligned}
\end{equation}
All these are operator paths on $E$. The covariance mean is not a pointwise
original potential. Write $U_D$ for left Young evolution and $P_D(t,s) =
S_E(t)U_D(t,s)S_E(-s)$.

\LeanMarker{L:symplectic-free}
\begin{lemma}[Free wave form and compact Young baseline]
  \label{L:symplectic-free}
  There is a bounded isometric skew-adjoint $J_E$ on $E$, with $J_E^2 = -I$,
  such that
  \[
    \omega(z,w)=\langle z,J_Ew\rangle_E
    =\langle u,g\rangle_{H^{1/2},H^{-1/2}}
    -\langle f,v\rangle_{H^{1/2},H^{-1/2}},
    \quad z=(u,v),\quad w=(f,g).
  \]
  The free group preserves $\omega$. For bounded real smooth $b$, the force
  $B(b)(u,v) = (0,bu)$ obeys $\omega(B(b)z,w)+\omega(z,B(b)w) = 0$.

  Let $1/2 < \theta < 1$ and let $X \in C^\theta([0,T];\mathcal{K}(E))$ start at
  zero, with $U_X$ its left Young evolution. Then
  \begin{equation}\label{E:symplectic-compact-young-baseline}
    U_X(t,s)-I\in\mathcal{K}(E),\qquad
    \|S_E(t)U_X(t,s)S_E(-s)\|
    \ge\|S_E(t-s)\|\ge1,\qquad 0\le s\le t\le T .
  \end{equation}
  More generally, for finitely many such paths $X_j$ starting at zero on $[0,B]$
  and lengths $0 \le \ell_j \le B$, put $P_j = S_E(\ell_j)U_{X_j}(\ell_j,0)$.
  Then
  \begin{equation}\label{E:symplectic-compact-product-baseline}
    P_N\cdots P_0-S_E\Big(\sum_{j=0}^N\ell_j\Big)\in\mathcal{K}(E),
    \qquad
    \|P_N\cdots P_0\|
    \ge\left\|S_E\Big(\sum_{j=0}^N\ell_j\Big)\right\|\ge1 .
  \end{equation}
  These norm bounds require neither Gaussian nor Hamiltonian operator paths.
\end{lemma}

\begin{proof}
  We identify the symplectic form in balanced coordinates and check free-wave
  preservation. We then show that compact Young perturbations remain compact
  at the endpoint. Translation to spatial infinity gives the required norm
  lower bound, including finite products.


  \medskip
  \noindent\textit{Step 1.~}We identify the form and its preservation by the
  free group.
  The balanced map $C_{\rm bal}(u,v) = (\Lambda^{1/2}u,\Lambda^{-1/2}v)$ is an
  onto real isometry to $L^2\times L^2$. Define $J_{\rm symp}(q,p) = (p,-q)$ and
  $J_E = C_{\rm bal}^{-1}J_{\rm symp}C_{\rm bal}$. This gives all stated
  operator properties and $J_E(f,g) = (\Lambda^{-1}g,-\Lambda f)$. The two
  displayed dualities follow from the balanced $L^2$ pairing. No $L^2$ pairing
  of a general distributional velocity is taken.

  In balanced Fourier coordinates, the free matrix is
  \[
    W_t(\xi)=
    \begin{pmatrix}
      \cos(tk)&\langle\xi\rangle\sin(tk)/k\\
      -k\sin(tk)/\langle\xi\rangle&\cos(tk)
    \end{pmatrix},\qquad k=|\xi|.
  \]
  At zero it is $\left(\begin{smallmatrix}1&t\\0&1\end{smallmatrix}\right)$, by
  the continuous sine quotient. Direct matrix multiplication gives
  $W_t^{\mathsf{T}}J_{\rm symp}W_t = J_{\rm symp}$, since its determinant is
  $\cos^2(tk)+\sin^2(tk) = 1$, also at zero. The entries are real, even and
  bounded on finite time horizons. Integrating the identity against real-phase
  Fourier vectors with Hermitian symmetry is legitimate by $L^2$
  Cauchy--Schwarz. It gives $\omega(S_E(t)z,S_E(t)w) = \omega(z,w)$ on all $E$.
  Equivalently one may first use Schwartz phases and then density. No rough
  generator derivative or frequency-zero trace is used.


  \medskip
  \noindent\textit{Step 2.~}We verify force skewness and compactness of Young
  endpoint differences.
  Finally, $\|bu\|_{H^{-1/2}} \le \|bu\|_2 \le \|b\|_\infty\|u\|_{H^{1/2}}$.
  Duality for these $L^2$ functions gives the absolutely integrable cancellation
  $-\int fbu+\int ubf = 0$. This is skewness for $\omega$, not for the $E$
  Hilbert inner product. The Young construction of
  Proposition~\ref{P:constant-young-linear} applies on $\mathcal{L}(E)$: its
  proof uses completeness and operator multiplication of norm one, exactly as on
  $\mathcal{L}(F)$. For fixed $s,t$, each left Young sum
  \[
    \sum_j\big(X(t_{j+1})-X(t_j)\big)U_X(t_j,s)
  \]
  is compact. The sums converge in operator norm to $U_X(t,s)-I$, so closedness
  of the compact ideal gives the first assertion
  in~\eqref{E:symplectic-compact-young-baseline}. Bounded composition then shows
  that the physical endpoint differs from $S_E(t-s)$ by a compact operator.


  \medskip
  \noindent\textit{Step 3.~}We bound the norm of a compact perturbation by the
  free norm.
  We prove the required compact-perturbation bound for this whole-space phase.
  Translations $\tau_y$ on both entries are isometries on $E$ and commute with
  $S_E(h)$. For fixed $z,w \in E$, their translated pairing is the Fourier
  transform of an $L^1$ function: the two weighted Fourier products are
  integrable by Cauchy--Schwarz in the Bessel norms. Approximation of that
  function in $L^1$ by compact smooth functions, followed by integration by
  parts for the approximants, proves $\langle\tau_yz,w\rangle_E \to 0$ as $|y|
  \to \infty$. Thus, $\tau_yz$ tends weakly to zero. For compact $K$,
  $\|K\tau_yz\| \to 0$: otherwise compactness supplies a norm-convergent
  subsequence with nonzero norm, while bounded linear testing makes its limit
  zero. For unit $z$ it follows that
  \[
    \|S_E(h)+K\|
    \ge\liminf_{|y|\to\infty}\|(S_E(h)+K)\tau_yz\|
    \ge\|S_E(h)z\|.
  \]
  Since translations commute with the free group and preserve the norm,
  $\|S_E(h)\tau_yz\|=\|S_E(h)z\|$. The compact term tends to zero on this
  weakly vanishing family. The reverse triangle inequality therefore gives the
  second inequality in the display for each fixed unit vector, before the
  supremum over vectors is taken.
  Taking the supremum over $z$ gives $\|S_E(h)+K\| \ge \|S_E(h)\|$. The
  isometric velocity reversal $\mathcal{R}(u,v) = (u,-v)$ satisfies $S_E(-h) =
  \mathcal{R} S_E(h)\mathcal{R}$ by the full free Fourier matrix. The group law
  yields $1 \le \|S_E(-h)\|\|S_E(h)\| = \|S_E(h)\|^2$. This
  proves~\eqref{E:symplectic-compact-young-baseline}.


  \medskip
  \noindent\textit{Step 4.~}We retain the compact difference through finite
  ordered products.
  Finally, $P_j-S_E(\ell_j)$ is compact. Telescoping the finite ordered product
  against its free product gives a finite sum of bounded--compact--bounded
  products, hence a compact difference. Only the free factors are combined by
  the group law; no noisy factors are commuted. The same translation argument
  proves \eqref{E:symplectic-compact-product-baseline}. A zero length gives $P_j
  = I$. All assertions are at fixed finite horizons and products; no uniform
  compact approximation in a radius, horizon or number of factors is asserted.
  This proves Lemma~\ref{L:symplectic-free}.
\end{proof}

\begin{lemma}[Covariance means and source primitives]
  \label{L:symplectic-source-mean}
  The expectations in~\eqref{E:symplectic-combined-driver} are Bochner defined
  and
  \begin{equation}\label{E:symplectic-mean-convergence}
    \|D_\Phi^\rho-D_\Phi\|_{C^\theta}
    \le\|\Phi\|_{H_W}\|M^\rho-M\|_{L^2C^\theta}\longrightarrow0 .
  \end{equation}
  For each positive $\rho$, there is a deterministic real time-continuous,
  spatially smooth $h_\Phi^\rho$ such that
  \begin{equation}\label{E:symplectic-source-primitive}
    D_\Phi^\rho(t)
    =\int_0^t S_E(-r)B(m_Rh_\Phi^\rho(r))S_E(r)\ud r .
  \end{equation}
  This is the statewise primitive and defines a Lipschitz operator path. The
  combined source coefficient is $b^\rho = m_R(cY^\rho+h_\Phi^\rho)$.
\end{lemma}

\begin{proof}
  We bound the Bochner covariance means by path moments. At positive source
  widths, the feature representation gives a smooth deterministic coefficient,
  whose ordinary rows identify the covariance primitive.


  \medskip
  \noindent\textit{Step 1.~}We bound the covariance mean and its source error.
  Use~\eqref{E:source-driver-limit} for $M^\rho\to M$ in $L^2C^\theta$,
  and~\eqref{E:source-transfer-norm} with $C=1$
  from~\eqref{E:source-two-dominations} for
  $\|M^\rho\|_{L^2C^\theta}\le\|M\|_{L^2C^\theta}$.
  Separable Gaussian path moments give $\E\|ZM\|_{C^\theta} \le
  \|\Phi\|_{H_W}\|M\|_{L^2C^\theta}$, and the same holds for sources. Bochner
  Cauchy--Schwarz gives~\eqref{E:symplectic-mean-convergence}. Covariance
  contraction also gives the uniform source bound $\|D_\Phi^\rho\|_{C^\theta}
  \le \|\Phi\|_{H_W}\|M\|_{L^2C^\theta}$.


  \medskip
  \noindent\textit{Step 2.~}We identify the regularized deterministic
  coefficient.
  The construction preserving the given noise has an isometry $j_{\rm src}:H_W
  \to \mathcal{G}$ and an isonormal process $\mathcal{I}$ with
  $\mathcal{I}(j_{\rm src}h) = W(h)$. Its source random field is $Y^\rho(t,x) =
  \mathcal{I}(v_{\rho;t,x})$, with continuous $\mathcal{G}$-valued spatial
  derivatives at positive widths. Consequently,
  \[
    h_\Phi^\rho(t,x)=\E[W(\Phi)Y^\rho(t,x)]
    =\langle j_{\rm src}\Phi,v_{\rho;t,x}\rangle_{\mathcal{G}}
  \]
  has the claimed regularity. Bounded linear testing commutes with those
  derivatives. After multiplication by $m_R$ the coefficient has one compact
  support and bounded derivatives on $[0,T]$.


  \medskip
  \noindent\textit{Step 3.~}We identify its ordinary primitive with the source
  covariance mean.
  Proposition~\ref{P:source-ordinary-response} identifies $M^\rho$ with the
  ordinary source primitive; the approximation here extends smooth time tests
  to the indicator wave tests in~\eqref{E:source-transferred-row}, with both
  limits taken in the original covariance completion.
  On Schwartz states, the row of the right side
  of~\eqref{E:symplectic-source-primitive} is the integral of $h_\Phi^\rho$
  against the explicit masked source wave product. Ordinary Fubini is justified
  by $\E|ZY^\rho(t,x)| \le \|\Phi\|_{H_W} \sqrt{\E|Y^\rho(t,x)|^2}$ and bounded
  point variance on this fixed compact set at positive widths. The source test
  and indicator approximation extends this to the completed rows by Gaussian
  Cauchy--Schwarz and covariance contraction. Hence, this row is the row of
  $\E[ZM^\rho(t)]$. Countable dense states and rational times, then continuity,
  identify the operator path. The uniformly bounded integrand gives its
  Lipschitz bound. This identifies the combined source coefficient exactly. The
  combined source is generally noncentered and the mask occurs once. This proves
  Lemma~\ref{L:symplectic-source-mean}.
\end{proof}

\begin{proposition}[Source conservation]\label{P:symplectic-source}
  At each fixed positive $\rho$, on one probability-one set for the fixed
  $c,\Phi,R,T$, the physical source flow satisfies
  \[
    P_{D^\rho}(t,s)^*J_EP_{D^\rho}(t,s)=J_E,
    \qquad 0\le s\le t\le T .
  \]
\end{proposition}

\begin{proof}
  Intersect the source coefficient regularity event with the event identifying
  its first primitive with the ordinary primitive for all times and states. The
  source construction supplies both at this fixed $\rho$. On their intersection
  $b^\rho$ has a common compact support, is continuous in supremum norm and is
  uniformly bounded on $[0,T]$. Thus, $C_b(r) = S_E(-r)B(b^\rho(r))S_E(r)$ is
  uniformly bounded and strongly continuous on states.
  Lemma~\ref{L:symplectic-free} gives $\omega(C_b(r)z,w)+\omega(z,C_b(r)w) = 0$
  for all $r,z,w$.

  Free preservation transfers the force identity to its interaction conjugate:
  applying $S_E(r)$ to both arguments turns the two terms with $C_b(r)$ into
  the two skew-force terms with $B(b^\rho(r))$. Their sum is zero. This uses
  the symplectic form, not skew-adjointness in the phase Hilbert norm.
  The source and mean primitive identities identify the Young state equation
  with $U(t,s)e = e+\int_s^t C_b(r)U(r,s)e\ud r$. Indeed, replace each operator
  path increment in a Young sum by its statewise integral. The error is bounded
  by the uniform coefficient bound times the uniform continuity error of $U$ on
  the partition. Each state trajectory is therefore $C^1$, with derivative
  $C_b(t)U(t,s)e$, by the strong-integral difference quotient. Differentiating
  the bounded bilinear form $\omega(U(t,s)e,U(t,s)f)$ gives zero. This holds on
  the same path set for all states and time pairs, since the bounded ordinary
  evolution is defined there on all $E$. Free conjugation gives the displayed
  physical identity. No uncountable intersection of state-dependent null sets or
  derivative of the free generator on a rough state was used. This proves
  Proposition~\ref{P:symplectic-source}.
\end{proof}

Proposition~\ref{P:symplectic-source} concerns the regularized flow
$P_{D^\rho}$ at a fixed positive $\rho$. We now prove the same conservation
identity for the flow $P_D$ itself, together with a lower bound on its norm.

\begin{proposition}[Conservation, inverse and a lower norm bound]
  \label{P:symplectic-raw}
  For the fixed parameters in~\eqref{E:symplectic-combined-driver}, on one
  probability-one set,
  \begin{equation}\label{E:symplectic-raw-conservation}
    P_D(t,s)^*J_EP_D(t,s)=J_E,\qquad
    \|P_D(t,s)\|\ge1,\qquad 0\le s\le t\le T .
  \end{equation}
  The flow is invertible. For $c = 0$ these assertions are deterministic for
  every fixed $\Phi$.
\end{proposition}

\begin{proof}
  We first show that every driver increment satisfies the same closed linear
  symplectic constraint. We pass that constraint to the Gaussian path and its
  covariance mean, then prove conservation by summing Young remainder
  estimates. Invertibility comes separately from the Young construction.


  \medskip
  \noindent\textit{Step 1.~}We identify the closed constraint for the
  coefficient paths.
  We first identify a closed real linear constraint on the operator path. Write
  \[
    \mathfrak{g}_J=\{A\in\mathcal{L}(E):A^*J_E+J_EA=0\}.
  \]
  The coefficient $H_{j,R}$ in~\eqref{E:driver-covariance-coefficient} is real
  and smooth. Free conservation gives $S_E(-r)^*J_E = J_ES_E(r)$ and $J_ES_E(-r)
  = S_E(r)^*J_E$. Lemma~\ref{L:symplectic-free} therefore makes
  $S_E(-r)B(H_{j,R}(r))S_E(r)$ belong to $\mathfrak{g}_J$. Testing the statewise
  integral in~\eqref{E:driver-coefficient-operator} against two states shows
  $A_j(0,t) \in \mathfrak{g}_J$ for every $j,t$. No operator-norm derivative of
  the free group is taken. The real finite sums
  in~\eqref{E:driver-path-convergence} obey the same constraint. The map
  $A\mapsto A^*J_E+J_EA$ is bounded on the path space, so their $L^2C^\theta$
  limit $M$ obeys it at all times on one probability-one path set.


  \medskip
  \noindent\textit{Step 2.~}We pass the constraint to the combined driver.
  The mean is a Bochner expectation since $\E\|W(\Phi)M\|_{C^\theta} \le
  \|\Phi\|_{H_W}\|M\|_{L^2C^\theta} < \infty$. Bounded real linear maps commute
  with this expectation, hence $D_\Phi(t) \in \mathfrak{g}_J$ for every $t$.
  Thus, $D = cM+D_\Phi$ starts at zero and is a $C^\theta$ path in
  $\mathfrak{g}_J$ on that path set. Complex Gaussian coefficients are not used.
  All these variables remain on the given $W$.


  \medskip
  \noindent\textit{Step 3.~}We prove conservation from the Young increment
  expansion.
  The proof of Proposition~\ref{P:constant-young-linear} applies on
  $\mathcal{L}(E)$, with multiplication of norm one. Fix any $s$ and put $V(r) =
  U_D(r,s)$ for $r \ge s$. The Young remainder on $[u,v] \subset [s,T]$ gives
  \[
    V(v)-V(u)=D_{u,v}V(u)+R_{u,v},\qquad
    \|R_{u,v}\|\le C|v-u|^{2\theta}.
  \]
  The exact increment of $V^*J_EV$ has first-order term
  $V(u)^*(D_{u,v}^*J_E+J_ED_{u,v})V(u) = 0$. Its remaining terms have norm at
  most $C'|v-u|^{2\theta}$, using the bounded $C^\theta$ norm of $V$. Their sum
  on a partition is at most $C'(t-s)\operatorname{mesh}^{2\theta-1}$, which
  tends to zero. Therefore, $U_D(t,s)^*J_EU_D(t,s) = J_E$. This deterministic
  argument applies to every time pair on the same operator path; it requires no
  state-dependent null sets. Free conservation in Lemma~\ref{L:symplectic-free}
  gives \eqref{E:symplectic-raw-conservation}, including $1 = \|J_E\| \le
  \|P_D(t,s)\|^2$.


  \medskip
  \noindent\textit{Step 4.~}We retain the independent justification of
  invertibility.
  Invertibility comes from the local Young Neumann pieces and their finite
  products, followed by the free inverse; it is not inferred from conservation
  in infinite dimension. For $c = 0$ the path and the argument are deterministic
  for each $\Phi$. Zero covariance gives the free flow directly. At each fixed
  $\delta$ one may take the given reference $W\circ T_{\delta,0}$, $R =
  \delta^{-1}+2$, $T = \delta^{-1}$, $c = g_\delta$ and $\Phi = 0$. No
  uniform-radius or uniform bound for controls assertion follows. This proves
  Proposition~\ref{P:symplectic-raw}.
\end{proof}

\begin{proposition}[The exact completed block mean]\label{P:symplectic-block-mean}
  Use the true partition $B,\delta,N,s_m,\ell_m$ of
  Section~\ref{S:actual-cameron-mean}, with integer $B \ge 1$. Let $i_I:H_I \to
  H_W$ be the supported closed inclusion and $T_m = T_{\delta,m}:H_{\rm fast}
  \to H_W$ the original isometry; $H_{\rm fast}$ has the same original
  covariance in fast coordinates. For $\Phi \in H_I$ set
  \[
    W_m(f)=W(T_mf),\qquad
    \phi_m=\delta^{-a/2}T_m^*i_I\Phi,\qquad
    \mathcal{T}_E\Phi=i^*T\Phi\,i ,
  \]
  where $T$ is the full $F$ mean of Proposition~\ref{P:cameron-full-mean}.
  Define
  \begin{equation}\label{E:symplectic-block-primitive}
    H_{m,E}(\Phi)(u)=\int_0^{\min(u,\ell_m)}
    S_E(-v)\mathcal{T}_E\Phi(s_m+\delta v)S_E(v)\ud v .
  \end{equation}
  With $M_m = M_R(W_m)$ and $D_{\phi_m}(u) = \E[M_m(u)W_m(\phi_m)]$, one has
  \begin{equation}\label{E:symplectic-block-direction}
    D_{\phi_m}(\min(u,\ell_m))=H_{m,E}(\Phi)(u),\qquad 0\le u\le B .
  \end{equation}
  Consequently, the pure-mean endpoint obeys
  \[
    \|S_E(\ell_m)U_{H_{m,E}(\Phi)}(\ell_m,0)\|\ge1 .
  \]
\end{proposition}

\begin{proof}
  We define the fast covariance direction using the adjoint of the original
  isometry. We identify the stopped completed mean by bounded compression,
  then remove the orthogonal complement of the fast first chaos from its
  covariance pairing.


  \medskip
  \noindent\textit{Step 1.~}We define the fast control and its completed mean
  path.
  The original pullback~\eqref{E:raw-block-pullback} has covariance factor
  $\delta^{a-2}\delta^2\delta^{-a} = 1$, hence respects nulls and extends
  isometrically from fast to slow original space. Its Hilbert adjoint is a
  contraction. Thus, $\|\phi_m\| \le \delta^{-a/2}\|\Phi\|$ is finite at fixed
  $\delta$, without a uniform control-radius assertion. The fast noise remains
  on the same given $W$. For applying the preceding propositions, construct
  regularized noises with $W_m$ as the given reference, extending by an
  independent auxiliary source while retaining all base variables. No
  slow-source pullback is presumed; the original covariance direction is
  independent of the auxiliary realization.

  The compact-conjugation and primitive estimates are those in
  Proposition~\ref{P:cameron-actual-embedding},
  including~\eqref{E:cameron-local-primitive-bound}, restricted from $F$ to
  $E$: compactness gives norm-continuous conjugation and Hölder integration
  gives exponent $1-1/r$.
  The mean $T\Phi$ has range in closed $iE$. By
  Proposition~\ref{P:cameron-full-mean}, bounded compression passes its strong
  measurability, $L^r$ bound and compact operator values to $\mathcal{T}_E\Phi
  \in L^r(I;\mathcal{K}(E))$. Choose the full parameter
  choices~\eqref{E:cameron-block-parameters}. Its intervals are nonempty because
  $\theta < 1-a/2$ and $\theta > 1/2$; $p$ only instantiates that package here.
  The compact-conjugation argument makes \eqref{E:symplectic-block-primitive} a
  Bochner integral, and the $L^r$ primitive estimate gives $C^\theta$, since
  $1-1/r > \theta$. Affine pullback preserves null sets. Its integrand is
  compact valued almost everywhere; closedness of $\mathcal{K}(E)$ under Bochner
  integration makes $H_{m,E}(\Phi)$ a compact-valued path starting at zero, also
  after constant stopping.


  \medskip
  \noindent\textit{Step 2.~}We compute the unscaled covariance by compressing
  the full phase.
  Here $A_{F,m}$ is the full $F$ interaction operator path from
  Proposition~\ref{P:constant-phase-driver}, and $H_{m,F}$ is its stopped
  completed mean primitive in~\eqref{E:cameron-actual-covariance}. The completed
  identity~\eqref{E:cameron-actual-covariance} at original projection rank zero
  gives, for $0 \le u \le \ell_m$,
  \[
    \E[g_\delta A_{F,m}(u)W(i_I\Phi)]
    =\sqrt\delta\,H_{m,F}(\Phi)(u).
  \]
  Compress through $i$ using the operator path intertwining, $S_Fi = iS_E$, and
  the mean range in $iE$ preserved by $S_F$. Bochner expectation commutes with
  this bounded compression. The constant columns are retained in the full $F$
  operator path and mean until this compression; they are inactive only on the
  $iE$ inputs. Since $g_\delta\delta^{a/2} = \sqrt\delta$ and $g_\delta > 0$,
  division gives the unscaled covariance
  \begin{equation}\label{E:symplectic-unscaled-covariance}
    \E[M_m(u)W(i_I\Phi)]
    =\delta^{a/2}H_{m,E}(\Phi)(u),\qquad u\le\ell_m .
  \end{equation}
  The factor here is $\delta^{a/2}$, not $\sqrt\delta$. Completion in the
  earlier covariance identity permits every $\Phi \in H_I$, without positivity
  or smoothness.


  \medskip
  \noindent\textit{Step 3.~}We identify the covariance direction in the fast
  noise space.
  The operator $T_mT_m^*$ is the orthogonal projection onto the original
  isometry's closed range. Thus, $W((I-T_mT_m^*)i_I\Phi)$ is orthogonal to the
  entire fast first chaos. Every continuous linear coordinate of $M_m$ is an
  $L^2$ limit of finite $W_m$-basis sums. Its covariance with this difference
  vanishes. Bochner covariance exists by $L^2$ path moments, and bounded linear
  functionals separate its values. Hence,
  \[
    \E[M_m(u)W(i_I\Phi)]
    =\E[M_m(u)W_m(T_m^*i_I\Phi)] .
  \]
  Multiply by $\delta^{-a/2}$ and use \eqref{E:symplectic-unscaled-covariance}
  to obtain \eqref{E:symplectic-block-direction} on $[0,\ell_m]$, then stop both
  sides. No unstopped equality beyond $\ell_m$ is asserted.


  \medskip
  \noindent\textit{Step 4.~}We apply the compact endpoint lower bound.
  Apply Lemma~\ref{L:symplectic-free} directly to the compact-valued path
  $H_{m,E}(\Phi)$ starting at zero on $[0,B]$. Its physical endpoint has norm at
  least $\|S_E(\ell_m)\| \ge 1$. If $\ell_m = 0$, the endpoint is the identity
  and its norm is one. The radius may be chosen $\delta^{-1}+2$ after fixing
  $\delta$; all identities are fixed-parameter instances. This proves
  Proposition~\ref{P:symplectic-block-mean}.
\end{proof}

For zero covariance the operator paths and covariance means vanish, and the
claims reduce directly to the free identity in Lemma~\ref{L:symplectic-free}.
These results give a lower bound on the operator norm and its exact completed
mean application. They provide no uniform bound on the control radius,
stochastic moment upper bound or nonlinear product ratio by themselves.
